\documentclass[a4paper,11pt]{article}
\usepackage[a4paper,margin=2cm]{geometry}
\usepackage[T1]{fontenc}
\usepackage[utf8]{inputenc}
\usepackage[ngerman,english]{babel}
\usepackage{lmodern}
\usepackage{microtype}
\usepackage{xcolor}
\usepackage{parskip}
\usepackage{enumitem}
\usepackage[most]{tcolorbox}
\usepackage{titlesec}
\usepackage[hidelinks]{hyperref}
\definecolor{HeaderBlueDark}{RGB}{70,110,170}
\definecolor{RowBlueLight}{RGB}{235,243,255}
\definecolor{RowBlueDark}{RGB}{220,233,250}
\definecolor{SoftGray}{RGB}{90,90,90}
\setlength{\parindent}{0pt}
\setlist[itemize]{leftmargin=1.4em,itemsep=0.35em,topsep=0.35em}
\linespread{1.06}
\titleformat{\section}[block]{\Large\bfseries\color{HeaderBlueDark}}{}{0pt}{}
\titlespacing{\section}{0pt}{1.3em}{0.8em}
\newtcolorbox{exbox}{enhanced,breakable,colback=RowBlueLight,colframe=RowBlueDark,boxrule=0.4pt,arc=1.5mm,left=1.2mm,right=1.2mm,top=1mm,bottom=1mm,title={Beispiele \textnormal{(Examples)}},fonttitle=\bfseries,colbacktitle=RowBlueDark,coltitle=black}
\NewDocumentEnvironment{grammarentry}{mm+b}{%
  \begin{tcolorbox}[enhanced,breakable,colback=white,colframe=HeaderBlueDark,boxrule=0.6pt,arc=1.5mm,left=1.5mm,right=1.5mm,top=1mm,bottom=1mm,colbacktitle=HeaderBlueDark,coltitle=white,fonttitle=\bfseries,title={#1\hfill\normalfont\footnotesize #2}]
  #3
  \end{tcolorbox}
}{}
\newcommand{\GermanText}[1]{\textbf{Deutsch: }#1\par}
\newcommand{\EnglishText}[1]{{\small\color{SoftGray}\textbf{English: }#1\par}}
\newcommand{\ExampleItem}[2]{\item \textbf{#1}\\{\small\color{SoftGray}#2}}
\title{Präpositionen: \glqq trotz\grqq{} \textnormal{(Prepositions: \glqq trotz\grqq)}}
\date{}
\begin{document}
\selectlanguage{ngerman}
\maketitle
\tableofcontents
\newpage

\section{Trotz: Grundfunktion \textnormal{(Trotz: basic function)}}
\begin{grammarentry}{\textbf{trotz}}{konzessive Präposition}
\GermanText{\textbf{Trotz} nennt einen Umstand, der normalerweise gegen ein Ergebnis spricht, obwohl dieses Ergebnis dennoch eintritt. Die Präposition bedeutet \glqq ungeachtet\grqq{} oder \glqq in spite of\grqq.}
\EnglishText{\textbf{Trotz} introduces a circumstance that would normally work against an outcome, even though that outcome still happens. It means \glqq in spite of\grqq.}
\end{grammarentry}

\section{Kasus \textnormal{(Case)}}
\begin{grammarentry}{Standard: Genitiv}{Umgangssprachlich auch Dativ}
\GermanText{In der Standardsprache steht \textbf{trotz} normalerweise mit dem Genitiv. In der gesprochenen Umgangssprache ist auch der Dativ verbreitet. Für formelles Schreiben und Prüfungen ist der Genitiv die sichere Wahl.}
\EnglishText{In standard German, \textbf{trotz} normally takes the genitive. The dative is also common in colloquial speech. The genitive is the safe choice in formal writing and exams.}
\end{grammarentry}
\begin{exbox}
\begin{itemize}
\ExampleItem{Trotz des Regens gingen wir spazieren.}{Despite the rain, we went for a walk.}
\ExampleItem{Trotz der hohen Preise ist das Café gut besucht.}{Despite the high prices, the café is busy.}
\ExampleItem{Umgangssprachlich: trotz dem Regen.}{Colloquial: despite the rain.}
\end{itemize}
\end{exbox}

\section{Unterschied zu \glqq obwohl\grqq{} und \glqq trotzdem\grqq{} \textnormal{(Difference from \glqq obwohl\grqq{} and \glqq trotzdem\grqq)}}
\begin{grammarentry}{Ähnliche Bedeutung, unterschiedliche Grammatik}{Satzmuster beachten}
\GermanText{\textbf{Trotz} steht vor einer Nominalgruppe. \textbf{Obwohl} leitet einen Nebensatz ein; das konjugierte Verb steht am Ende. \textbf{Trotzdem} verbindet Hauptsätze; danach steht das konjugierte Verb auf Position 2.}
\EnglishText{\textbf{Trotz} precedes a noun phrase. \textbf{Obwohl} introduces a subordinate clause, where the finite verb goes to the end. \textbf{Trotzdem} connects main clauses; its finite verb comes in position two.}
\end{grammarentry}
\begin{exbox}
\begin{itemize}
\ExampleItem{Trotz des Regens gingen wir spazieren.}{Despite the rain, we went for a walk.}
\ExampleItem{Obwohl es regnete, gingen wir spazieren.}{Although it was raining, we went for a walk.}
\ExampleItem{Es regnete. Trotzdem gingen wir spazieren.}{It was raining. Nevertheless, we went for a walk.}
\end{itemize}
\end{exbox}

\section{Weiterführende Erklärung \textnormal{(Detailed guide)}}
\GermanText{Eine ausführliche Erklärung zu Bedeutung, Kasus, Adjektivendungen, Satzstellung, Kommas und festen Ausdrücken findest du in der Datei \href{https://github.com/donkarlo/nd_german_learning_project/blob/main/src/part/grammar/part/praeposition/kind/konzessiv/trotz/trotz.tex}{\textbf{trotz.tex}}.}
\EnglishText{A detailed guide to meaning, case, adjective endings, word order, commas, and fixed expressions is available in \href{https://github.com/donkarlo/nd_german_learning_project/blob/main/src/part/grammar/part/praeposition/kind/konzessiv/trotz/trotz.tex}{\textbf{trotz.tex}}.}

\section{Zusammenfassung \textnormal{(Summary)}}
\begin{grammarentry}{Merksatz}{kurz und wichtig}
\GermanText{\textbf{trotz + Genitiv + Nomen}; \textbf{obwohl + Nebensatz mit Verb am Ende}; \textbf{trotzdem + Hauptsatz mit Verb auf Position 2}.}
\EnglishText{\textbf{trotz + genitive + noun}; \textbf{obwohl + subordinate clause with the verb at the end}; \textbf{trotzdem + main clause with the verb in position two}.}
\end{grammarentry}
\end{document}
